% !TEX root = main.tex
% =============================================================================
% Introduction / background
% =============================================================================

\section{Introduction}
\label{sec:intro}

Additive manufacturing (AM) allows compact hydraulic components whose internal
channels follow free-form three-dimensional paths rather than intersecting
drilled holes.
A representative design case is multi-group channel layout in rotary
direct-drive servo valves: terminals are partitioned into groups; pipes of
different groups must remain geometrically separated for all poses of interest,
while pipes within a group may share trunks and branch into tree-like
manifolds.
Once a channel is placed it permanently occupies volume, so later groups must
clear earlier geometry.

Our companion work on discrete multi-group pipe routing under permanent
geometric occupancy~\cite{deng2025pipe_design} addresses the combinatorial
layout problem: given ports, obstacles, and group topologies, produce a
conflict-free coarse embedding on a constructive 3D grid benchmark, evaluated
by fixed-budget success of CBS, PBS, and priority planning.
That paper explicitly leaves post-layout smoothing as orthogonal follow-on
work.
The layouts it returns are feasible but still grid-angular; the present report
refines those centerlines.

\paragraph{Why robotic path smoothing is not enough.}
Trajectory fairing methods such as CHOMP, STOMP, TEB, and TrajOpt typically
optimize a single open curve for a temporary agent under time-indexed or soft
obstacle costs~\cite{zucker2013chomp,kalakrishnan2011stomp,roesmann2013teb,schulman2014trajopt}.
Pipe refinement differs along several axes that we make precise in
Sec.~\ref{sec:problem}: permanent tubular occupancy, branched topology with
degree-$\ge 3$ junctions, costs driven by declared usages rather than one
$s$--$t$ path.
The central geometric conflict is not ``make one spline smoother,'' but
\emph{how several overlapping usages share a trunk while remaining fair at every
declared terminal-to-terminal flow}.

\paragraph{Scope of this report.}
We formalize topology-preserving, multi-usage geometric refinement
under permanent tubular occupancy (Sec.~\ref{sec:problem}), survey adjacent
literatures (Sec.~\ref{sec:related}), and describe a hybrid trunk-share /
junction-split Hermite parameterization (Sec.~\ref{sec:method}) with soft path
and clearance costs (Sec.~\ref{sec:costs}).
Gradient-based experiments on synthetic instances with capsule clearance follow
in later sections.
We do not claim a new optimal continuous planner, commercial CAD parity, or
CFD-validated hydraulic improvement; those belong to later CAD and fluid
optimization, not to this geometric pass.

\paragraph{Contributions.}
\begin{enumerate}
  \item \textbf{Problem.} A clear formulation of post-routing centerline
        refinement for branched, multi-usage, permanently occupied tubes,
        including the shared-junction versus locally split junction regimes.
  \item \textbf{Method sketch.} Hybrid trunk-share / junction-split variables
        with soft length/curvature and clearance costs, hard ports and
        topology by construction, and hard clearance at accept
        (detailed in later sections).
  \item \textbf{Empirical note.} On hard multi-group synthetic cases, gradient
        refine with capsule clearance yields feasible designs with
        $\kappa r_{\max}\lesssim 1$ and large relative objective drop
        (Sec.~\ref{sec:experiments}).
\end{enumerate}
